% =====================================================================
%  Thème 5 — Résolution des triangles — NIVEAU DIFFICILE
%  Corrigé
% =====================================================================
\documentclass[11pt,a4paper]{article}
\usepackage[utf8]{inputenc}
\usepackage[T1]{fontenc}
\usepackage{lmodern}
\usepackage[french]{babel}
\usepackage{amsmath,amssymb,mathtools}
\usepackage{icomma}
\usepackage{xcolor}
\usepackage{colortbl}
\usepackage{tikz}
\usepackage[most]{tcolorbox}
\usepackage{enumitem}
\usepackage{array}
\usepackage[a4paper,margin=2.2cm]{geometry}
\usetikzlibrary{arrows.meta,calc}

\definecolor{primary}{RGB}{214,51,108}
\definecolor{primarydark}{RGB}{140,20,64}
\definecolor{excol}{RGB}{0,135,90}

\DeclareMathOperator{\tg}{tg}
\DeclareMathOperator{\cotg}{cotg}
\newcommand{\degr}{^{\circ}}
\newcommand{\Z}{\mathbb{Z}}

\newcounter{exo}
\newcommand{\exo}[1]{\medskip\refstepcounter{exo}%
  \noindent{\large\bfseries\color{primarydark}Exercice \theexo}%
  \ \textcolor{primary}{\textbullet}\ \textbf{#1}\par\nopagebreak\smallskip}
\newcommand{\rep}{\textcolor{excol}{$\blacktriangleright$}\ }

\setlength{\parindent}{0pt}
\pagestyle{empty}

\begin{document}

\begin{center}
{\LARGE\bfseries\color{primary}Thème 5 — Résolution des triangles}\\[2pt]
{\color{primarydark}\textsc{Niveau Difficile — Corrigé détaillé}}
\end{center}
\hrule height 1pt
\medskip

\exo{Un terrain triangulaire}
\begin{itemize}[leftmargin=1.4em,itemsep=2pt]
\item[\textbf{(a)}] \rep $\overline{BC}$ est opposé à $\hat A$, encadré par $\overline{AB}$ et
   $\overline{AC}$. Avec $\cos120\degr=-\tfrac12$ :
   \[
   \overline{BC}^2=6^2+6^2-2\cdot6\cdot6\cos120\degr=72-72\cdot(-\tfrac12)=72+36=108
   \ \Longrightarrow\ \overline{BC}=6\sqrt3.
   \]
\item[\textbf{(b)}] \rep $\mathcal{A}=\dfrac12\,\overline{AB}\cdot\overline{AC}\cdot\sin\hat A
   =\dfrac12\cdot6\cdot6\cdot\sin120\degr=18\cdot\dfrac{\sqrt3}{2}=9\sqrt3$.
\item[\textbf{(c)}] \rep Le triangle étant isocèle en $A$, $\hat B=\hat C$, et
   $\hat B+\hat C=180\degr-120\degr=60\degr$, donc $\hat B=\hat C=30\degr$.
   Vérification (loi des sinus) : $\dfrac{\sin\hat B}{\overline{AC}}=\dfrac{\sin30\degr}{6}=\dfrac1{12}$
   et $\dfrac{\sin\hat A}{\overline{BC}}=\dfrac{\sin120\degr}{6\sqrt3}=\dfrac{\sqrt3/2}{6\sqrt3}=\dfrac1{12}$ : cohérent.
\item[\textbf{(d)}] \rep $\mathcal{A}=\dfrac12\,\overline{BC}\cdot h$, donc
   $h=\dfrac{2\mathcal{A}}{\overline{BC}}=\dfrac{2\cdot9\sqrt3}{6\sqrt3}=\dfrac{18\sqrt3}{6\sqrt3}=3$
   (soit $30$ m).
\end{itemize}

\exo{La largeur d'une rivière}
\begin{itemize}[leftmargin=1.4em,itemsep=2pt]
\item[\textbf{(a)}] \rep $\hat C=180\degr-45\degr-45\degr=90\degr$.
\item[\textbf{(b)}] \rep $\overline{AB}=100$ est opposé à $\hat C=90\degr$. Par la loi des sinus,
   $\overline{AC}$ (opposé à $\hat B=45\degr$) et $\overline{BC}$ (opposé à $\hat A=45\degr$) valent
   \[
   \overline{AC}=\overline{BC}=100\cdot\frac{\sin45\degr}{\sin90\degr}
   =100\cdot\frac{\sqrt2/2}{1}=50\sqrt2\approx70{,}7\text{ m}.
   \]
\item[\textbf{(c)}] \rep $\overline{AC}^2+\overline{BC}^2=(50\sqrt2)^2+(50\sqrt2)^2=5000+5000=10000=100^2
   =\overline{AB}^2$ : le triangle est bien rectangle en $C$.
\item[\textbf{(d)}] \rep Aire $=\dfrac12\,\overline{AC}\cdot\overline{BC}\cdot\sin\hat C
   =\dfrac12\cdot50\sqrt2\cdot50\sqrt2\cdot1=\dfrac12\cdot5000=2500$ m$^2$.
   Or aire $=\dfrac12\,\overline{AB}\cdot(\text{largeur})$, d'où
   largeur $=\dfrac{2\cdot2500}{100}=50$ m.
\end{itemize}

\exo{Résolution complète d'un triangle}
\begin{itemize}[leftmargin=1.4em,itemsep=2pt]
\item[\textbf{(a)}] \rep Les côtés encadrant $\hat C$ sont $a$ et $b$ :
   \[
   \mathcal{A}=\frac12ab\sin\hat C\Rightarrow 18\sqrt3=\frac12\cdot12\cdot b\cdot\frac{\sqrt3}{2}
   =3b\sqrt3\ \Longrightarrow\ b=6.
   \]
\item[\textbf{(b)}] \rep $c=\overline{AB}$ opposé à $\hat C$ :
   \[
   c^2=a^2+b^2-2ab\cos60\degr=144+36-2\cdot72\cdot\tfrac12=180-72=108\ \Longrightarrow\ c=6\sqrt3.
   \]
\item[\textbf{(c)}] \rep Loi des sinus : $\dfrac{\sin\hat B}{b}=\dfrac{\sin\hat C}{c}$, donc
   \[
   \sin\hat B=b\cdot\frac{\sin\hat C}{c}=6\cdot\frac{\sin60\degr}{6\sqrt3}
   =6\cdot\frac{\sqrt3/2}{6\sqrt3}=\frac12\ \Longrightarrow\ \hat B=30\degr.
   \]
   Puis $\hat A=180\degr-60\degr-30\degr=90\degr$.
\item[\textbf{(d)}] \rep Comme $\hat A=90\degr$, le triangle est \textbf{rectangle en $A$}.
   Vérification par Pythagore ($a=\overline{BC}=12$ est l'hypoténuse) :
   $b^2+c^2=36+108=144=12^2=a^2$. Enfin, la hauteur issue de $A$ vérifie
   $\mathcal{A}=\dfrac12\,a\cdot h$, donc $h=\dfrac{2\cdot18\sqrt3}{12}=3\sqrt3\approx5{,}20$.
\end{itemize}

\end{document}
